% GENERATED FILE. Edit canonical lessons, then run npm run generate:books.
% canonical-source-hash: fnv1a64:edb8dc58c311b5c6
% canonical-lessons: TA-C447-munpanam, TA-C447-tavanai, TA-C447-pakki, TA-C447-miti, TA-C447-tokai

\chapter{Paying in Parts}
\label{ch:ta-paying-in-parts}
\hlchaptermodality{447}

\begin{chapteropening}
\textbf{By the end of this chapter:} \emph{I can talk about an advance, instalments, dues, the remainder and the total sum.}

Complete the last lesson of chapter 447: I can talk about an advance, instalments, dues, the remainder and the total sum.
\end{chapteropening}

\section[\texorpdfstring{\ta{முன்பணம்} (muṉpaṇam)}{முன்பணம் (muṉpaṇam)}]{\ta{முன்பணம்} (muṉpaṇam) — an advance payment; a deposit}

\label{lesson:TA-C447-munpanam}



\begin{quote}
Before the new one: say the Tamil for a cheque, then the Tamil for cash.
\end{quote}

\subsection*{You'll want to know: \ta{முன்பணம்}}
\textbf{\ta{முன்பணம்}} — \emph{muṉpaṇam} — \textquotedblleft{}an advance payment; a deposit\textquotedblright{}.

\begin{quote}
\ta{அறைக்கு} \ta{ஆயிரம்} \ta{ரூபாய்} \ta{முன்பணம்} \ta{கொடுத்தோம்}.
\end{quote}

\begin{quote}
\emph{aṟaikku āyiram rūpāy muṉpaṇam koṭuttōm.}
\end{quote}

\begin{quote}
We paid a thousand rupees in advance for the room.
\end{quote}

\ta{முன்}, before, + \ta{பணம்}, money.

\begin{cousinweb}
\textbf{\ta{முன்பணம்}} (\emph{muṉpaṇam}) is built from \ta{முன்}, \textquotedblleft{}before\textquotedblright{} (DEDR 5020), and \ta{பணம்}, \textquotedblleft{}money\textquotedblright{}.
\end{cousinweb}

\begin{grammarlens}[title={What you've built}]
One more word for getting around.
\end{grammarlens}

\subsection*{Your turn}
\begin{itemize}
\raggedright
  \item \emph{Say it:} \emph{muṉpaṇam}
  \item \emph{Say it:} \emph{muṉpaṇam}, once more
  \item \emph{Say it:} \emph{aṟaikku āyiram rūpāy muṉpaṇam koṭuttōm.} — We paid a thousand rupees in advance for the room.
  \item \emph{Recall:} say the Tamil for a cheque, then the Tamil for cash, then say \emph{muṉpaṇam} again
\end{itemize}

\begin{tcolorbox}[breakable,colback=teal!4,colframe=teal!35!black,title={Before you move on}]
What does \ta{முன்பணம்} mean? (\textquotedblleft{}An advance payment; a deposit\textquotedblright{}.) Say it once more.
\end{tcolorbox}
\section[\texorpdfstring{\ta{தவணை} (tavaṇai)}{தவணை (tavaṇai)}]{\ta{தவணை} (tavaṇai) — an instalment; a fixed term}

\label{lesson:TA-C447-tavanai}



\begin{quote}
Before the new one: say the Tamil for cash, then the Tamil for an advance payment.
\end{quote}

\subsection*{You'll want to know: \ta{தவணை}}
\textbf{\ta{தவணை}} — \emph{tavaṇai} — \textquotedblleft{}an instalment; a fixed term\textquotedblright{}.

\begin{quote}
\ta{நான்} \ta{கைப்பேசியை} \ta{மாதத்} \ta{தவணையில்} \ta{வாங்கினேன்}.
\end{quote}

\begin{quote}
\emph{nāṉ kaippēciyai mātat tavaṇaiyil vāṅkiṉēṉ.}
\end{quote}

\begin{quote}
I bought the mobile phone in monthly instalments.
\end{quote}

\ta{மாதத்} \ta{தவணை} is a monthly instalment.

\begin{cousinweb}
\textbf{\ta{தவணை}} (\emph{tavaṇai}) is native Dravidian (DEDR 3108), \textquotedblleft{}a fixed term for payment, an instalment\textquotedblright{}.
\end{cousinweb}

\begin{grammarlens}[title={What you've built}]
One more word for getting around.
\end{grammarlens}

\subsection*{Your turn}
\begin{itemize}
\raggedright
  \item \emph{Say it:} \emph{tavaṇai}
  \item \emph{Say it:} \emph{tavaṇai}, once more
  \item \emph{Say it:} \emph{nāṉ kaippēciyai mātat tavaṇaiyil vāṅkiṉēṉ.} — I bought the mobile phone in monthly instalments.
  \item \emph{Recall:} say the Tamil for cash, then the Tamil for an advance payment, then say \emph{tavaṇai} again
\end{itemize}

\begin{tcolorbox}[breakable,colback=teal!4,colframe=teal!35!black,title={Before you move on}]
What does \ta{தவணை} mean? (\textquotedblleft{}An instalment; a fixed term\textquotedblright{}.) Say it once more.
\end{tcolorbox}
\section[\texorpdfstring{\ta{பாக்கி} (pākki)}{பாக்கி (pākki)}]{\ta{பாக்கி} (pākki) — an amount due; arrears}

\label{lesson:TA-C447-pakki}



\begin{quote}
Before the new one: say the Tamil for an advance payment, then the Tamil for an instalment.
\end{quote}

\subsection*{You'll want to know: \ta{பாக்கி}}
\textbf{\ta{பாக்கி}} — \emph{pākki} — \textquotedblleft{}an amount due; arrears\textquotedblright{}.

\begin{quote}
\ta{கடைக்காரருக்கு} \ta{இன்னும்} \ta{நூறு} \ta{ரூபாய்} \ta{பாக்கி} \ta{இருக்கிறது}.
\end{quote}

\begin{quote}
\emph{kaṭaikkārarukku iṉṉum nūṟu rūpāy pākki irukkiṟatu.}
\end{quote}

\begin{quote}
A hundred rupees is still owed to the shopkeeper.
\end{quote}

\ta{பாக்கி} \ta{வை} is to leave a sum unpaid.

\begin{cousinweb}
\textbf{\ta{பாக்கி}} (\emph{pākki}) is a loan from Arabic \emph{bāqī}, \textquotedblleft{}remaining\textquotedblright{}, which reached Tamil through Hindustani.
\end{cousinweb}

\begin{grammarlens}[title={What you've built}]
One more word for getting around.
\end{grammarlens}

\subsection*{Your turn}
\begin{itemize}
\raggedright
  \item \emph{Say it:} \emph{pākki}
  \item \emph{Say it:} \emph{pākki}, once more
  \item \emph{Say it:} \emph{kaṭaikkārarukku iṉṉum nūṟu rūpāy pākki irukkiṟatu.} — A hundred rupees is still owed to the shopkeeper.
  \item \emph{Recall:} say the Tamil for an advance payment, then the Tamil for an instalment, then say \emph{pākki} again
\end{itemize}

\begin{tcolorbox}[breakable,colback=teal!4,colframe=teal!35!black,title={Before you move on}]
What does \ta{பாக்கி} mean? (\textquotedblleft{}An amount due; arrears\textquotedblright{}.) Say it once more.
\end{tcolorbox}
\section[\texorpdfstring{\ta{மீதி} (mīti)}{மீதி (mīti)}]{\ta{மீதி} (mīti) — the remainder; the rest}

\label{lesson:TA-C447-miti}



\begin{quote}
Before the new one: say the Tamil for an instalment, then the Tamil for an amount due.
\end{quote}

\subsection*{You'll want to know: \ta{மீதி}}
\textbf{\ta{மீதி}} — \emph{mīti} — \textquotedblleft{}the remainder; the rest\textquotedblright{}.

\begin{quote}
\ta{பாதி} \ta{பணம்} \ta{கொடுத்தேன்}, \ta{மீதியை} \ta{நாளை} \ta{கொடுப்பேன்}.
\end{quote}

\begin{quote}
\emph{pāti paṇam koṭuttēṉ, mītiyai nāḷai koṭuppēṉ.}
\end{quote}

\begin{quote}
I paid half the money; I will pay the rest tomorrow.
\end{quote}

\ta{மீதிச்} \ta{சில்லறை} is the change due.

\begin{cousinweb}
\textbf{\ta{மீதி}} (\emph{mīti}) is native Dravidian (DEDR 4838), \textquotedblleft{}the remainder, that which is left\textquotedblright{}.
\end{cousinweb}

\begin{grammarlens}[title={What you've built}]
One more word for getting around.
\end{grammarlens}

\subsection*{Your turn}
\begin{itemize}
\raggedright
  \item \emph{Say it:} \emph{mīti}
  \item \emph{Say it:} \emph{mīti}, once more
  \item \emph{Say it:} \emph{pāti paṇam koṭuttēṉ, mītiyai nāḷai koṭuppēṉ.} — I paid half the money; I will pay the rest tomorrow.
  \item \emph{Recall:} say the Tamil for an instalment, then the Tamil for an amount due, then say \emph{mīti} again
\end{itemize}

\begin{tcolorbox}[breakable,colback=teal!4,colframe=teal!35!black,title={Before you move on}]
What does \ta{மீதி} mean? (\textquotedblleft{}The remainder; the rest\textquotedblright{}.) Say it once more.
\end{tcolorbox}
\section[\texorpdfstring{\ta{தொகை} (tokai)}{தொகை (tokai)}]{\ta{தொகை} (tokai) — a sum; an amount}

\label{lesson:TA-C447-tokai}



\begin{quote}
Before the new one: say the Tamil for an amount due, then the Tamil for the remainder.
\end{quote}

\subsection*{You'll want to know: \ta{தொகை}}
\textbf{\ta{தொகை}} — \emph{tokai} — \textquotedblleft{}a sum; an amount\textquotedblright{}.

\begin{quote}
\ta{மொத்தத்} \ta{தொகை} \ta{எவ்வளவு} \ta{என்று} \ta{கேட்டேன்}.
\end{quote}

\begin{quote}
\emph{mottat tokai evvaḷavu eṉṟu kēṭṭēṉ.}
\end{quote}

\begin{quote}
I asked how much the total amount was.
\end{quote}

\ta{மொத்தத்} \ta{தொகை} is the total amount.

\begin{cousinweb}
\textbf{\ta{தொகை}} (\emph{tokai}) is native Dravidian (DEDR 3476), \textquotedblleft{}an assembly, a collection, a sum\textquotedblright{}.
\end{cousinweb}

\begin{grammarlens}[title={What you've built}]
One more word for getting around.
\end{grammarlens}

\subsection*{Your turn}
\begin{itemize}
\raggedright
  \item \emph{Say it:} \emph{tokai}
  \item \emph{Say it:} \emph{tokai}, once more
  \item \emph{Say it:} \emph{mottat tokai evvaḷavu eṉṟu kēṭṭēṉ.} — I asked how much the total amount was.
  \item \emph{Recall:} say the Tamil for an amount due, then the Tamil for the remainder, then say \emph{tokai} again
\end{itemize}

\begin{tcolorbox}[breakable,colback=teal!4,colframe=teal!35!black,title={Before you move on}]
What does \ta{தொகை} mean? (\textquotedblleft{}A sum; an amount\textquotedblright{}.) Say it once more.
\end{tcolorbox}
